\section{Plataforma experimental y antecedentes}
\subsection{Instrumentación}
La plataforma comprende un PSoC 5LP con ADC, DAC y temporizador de muestreo;
planta analógica con acondicionamiento de entrada y salida; fuente de
alimentación; PC con MATLAB; y osciloscopio \emph{Rigol DS1202Z-E}.
La instrumentación previa empleó USB-TMC/SCPI, sondas \(10\times\),
acoplamiento DC, telemetría UART a \SI{57600}{bit\per\second} y adquisición
sincronizada. Estos antecedentes constituyen una referencia operativa, no una
calibración vigente del nuevo montaje.

\subsection{Calibración histórica}
La Tabla~\ref{tab:calib_historica} reúne parámetros documentados en los
archivos de configuración y en el ajuste experimental del Laboratorio 5.
La adquisición del Lab7 comenzará con una \textbf{revalidación} de cero, rango,
pendientes y ruido a partir de su propio cableado.

\begin{table}[t]
\centering
\caption{Parámetros documentados del Laboratorio 5 (no revalidados en Lab7).}
\label{tab:calib_historica}
\small
\begin{tabular}{@{}ll@{}}
\toprule
Parámetro histórico & Valor documentado \\
\midrule
Offset del ADC & \SI{2.353802067}{\volt} \\
Escala ADC/planta & \SI{0.210251951}{\volt\per\volt} \\
Ganancia de DAC (+) & 3,842612 \\
Ganancia de DAC (-) & 2,640105691 \\
Offset físico J5 & \SI{17.209}{\milli\volt} \\
Pendiente J5/mando aplicado & 0,981710 \\
RMSE del ajuste de J7 & \SI{3.374}{\milli\volt} \\
\bottomrule
\end{tabular}
\end{table}

En las comparaciones se conservarán las tensiones crudas de los dos canales,
la telemetría digital, los valores calibrados, la referencia física, la
señal de control realmente aplicada y los metadatos de disparo del osciloscopio.
Esto evita confundir coincidencias visuales de amplitud con validaciones
cuantitativas de la respuesta del sistema.
